\chapter{Fundamentos Teóricos}

\section{Acústica de salas}

La Respuesta al Impulso de la Sala, denotada como $h(t)$, constituye una caracterización muy útil del comportamiento acústico de un recinto cuando se aproxima como un sistema Lineal Invariante en el Tiempo (LTI). Esta idealización no es estricta en todos los casos, pero resulta razonable mientras la geometría y las condiciones del recinto se mantienen estables. Bajo dicho marco aproximado, $h(t)$ representa la salida del sistema ante una excitación impulsiva ideal y, en consecuencia, encapsula tanto los mecanismos de propagación directa como los procesos de reflexión en fronteras materiales. Desde una perspectiva fenomenológica, la respuesta temporal se descompone en tres contribuciones: sonido directo, reflexiones tempranas (\textit{early reflections}) y reverberación tardía (\textit{late reverberation}).

El concepto de tiempo de mezcla, $t_{mix}$, establece el umbral físico a partir del cual la densidad temporal de reflexiones se vuelve suficientemente alta como para que el campo acústico transite hacia un régimen difuso y de naturaleza estocástica \cite{Kuttruff2016}. Antes de dicho instante, las reflexiones mantienen una estructura temporal relativamente resoluble y conservan información geométrica macroscópica del recinto; después de $t_{mix}$, la superposición densa de trayectorias reduce la interpretabilidad determinista individual de cada reflexión. En práctica de ingeniería acústica, este umbral se sitúa típicamente en el entorno de los primeros 50 ms, lo que respalda la viabilidad de inferir o extrapolar la cola reverberante a partir del segmento temprano de la RIR.

Dado que la reverberación tardía presenta un carácter estocástico, su análisis resulta más robusto cuando se formula en términos energéticos. En este contexto, se adopta la Curva de Decaimiento de Energía (EDC) según la formulación de Schroeder \cite{Schroeder1965}:

\begin{equation}
\mathrm{EDC}(t) = \int_t^\infty h^2(\tau) \, d\tau,
\end{equation}

lo que permite cuantificar la energía remanente de la respuesta al impulso a medida que avanza el tiempo. La representación en dB de la EDC facilita la estimación del tiempo de reverberación $T_{60}$, definido como el intervalo necesario para que la energía descienda 60 dB respecto al nivel inicial.

En recintos reales, la transición desde el régimen temprano hacia la reverberación tardía rara vez es perfectamente idealizada. Factores como la absorción frecuencial de los materiales, la dispersión de las superficies, la geometría del recinto y pequeñas variaciones en la temperatura o la humedad introducen desviaciones respecto al modelo LTI perfecto. No obstante, la aproximación LTI sigue siendo extremadamente útil porque permite describir la respuesta acústica con una función impulsional única y estable, suficientemente precisa para muchos fines de simulación y aprendizaje automático.

\subsection{Modelo de filtros}

Desde el punto de vista del procesamiento de señales, la RIR puede interpretarse como la respuesta impulsional de un sistema lineal que transforma una excitación sonora en una señal reverberante. Esta interpretación permite conectar la acústica de salas con herramientas clásicas como la convolución, las transformadas de Fourier y los bancos de filtros.

El modelo de filtrado supone que la señal observada puede descomponerse en contribuciones por bandas de frecuencia, cada una con un comportamiento de decaimiento diferente. Esta descomposición resulta especialmente apropiada para la reverberación tardía, donde las altas frecuencias tienden a disiparse antes que las bajas.

\subsection{Modelos de aprendizaje}

El modelado de audio en bruto es complejo debido a la alta dimensionalidad de las señales temporales y a su sensibilidad a las variaciones de fase. Para abordar esto, en el aprendizaje profundo se emplean espacios latentes que actúan como representaciones comprimidas de las características esenciales de los datos \cite{Goodfellow2016}. 

Bajo esta formulación, un codificador ($ \mathcal{E} $) mapea la señal de entrada $\mathbf{x}$ a un espacio compacto de menor dimensión:

\begin{equation}
\mathbf{z} = \mathcal{E}(\mathbf{x}),
\end{equation}

mientras que el decodificador ($ \mathcal{D} $) se encarga de reconstruir la señal a partir de dicha representación:

\begin{equation}
\hat{\mathbf{y}} = \mathcal{D}(\mathbf{z}).
\end{equation}

Este enfoque permite concentrar la información acústica más relevante, lo que facilita la extrapolación temporal y reduce el coste computacional del modelo.
